\RequirePackage[dvipsnames,svgnames]{xcolor}
\documentclass[12pt,a4paper,leqno]{report}

\usepackage[utf8]{inputenc}
\usepackage[T1]{fontenc}
\usepackage[framemethod=TikZ]{mdframed}
\usepackage{amsmath}
\usepackage{booktabs}
\usepackage{caption}
\usepackage[sc]{mathpazo}
\usepackage{titlesec}
\usepackage[top=2cm,left=2.5cm,right=2.5cm,bottom=2cm]{geometry}
\usepackage[english]{babel}
\usepackage[most]{tcolorbox}
\usepackage{setspace}
\usepackage{hyperref}
\hypersetup{
    unicode=true, colorlinks=true,
    linkcolor=MidnightBlue, citecolor=OliveGreen,
    urlcolor=cyan, filecolor=BrickRed,
    breaklinks=true, pdfborder={0 0 0}
}
\linespread{1.3}
\usepackage{ulem, eso-pic}
\usepackage{fancyhdr}

\pagestyle{fancy}
\renewcommand{\headrulewidth}{0pt}
\fancyhf{}
\fancyfoot[C]{\sffamily Page~\thepage}
\fancyfoot[LE,RO]{\sffamily Group 10}
\fancyhead[LO]{\scriptsize\color{black!80}\leftmark}
\fancyfoot[LO]{\copyright AIMS Senegal~|~2025~/~2026}
\fancypagestyle{plain}{\fancyhead{}\renewcommand{\headrulewidth}{0pt}}
\renewcommand{\footrulewidth}{0.5pt}

\usepackage{tikz}
\usetikzlibrary{shadows.blur}

\titleformat{\chapter}[block]{\LARGE\bfseries\sffamily}
{%
  \raisebox{4\height}{\sffamily\scriptsize\textcolor{violet}{\MakeUppercase{\chaptertitlename}}}\space
  \raisebox{-0.3\height}{%
    \tikz[baseline=(n.base)]{\node[circle,shading=axis,top color=violet,bottom color=brown,
    double,double distance=1.5pt,line width=0.7pt,inner sep=3pt,text=white,font=\mdseries\large,
    blur shadow={shadow xshift=0.7pt,shadow yshift=-0.7pt,shadow blur steps=3,shadow opacity=0.8}
    ](n){\bigchapternumber};}}\space
}{0pt}{\printtitleNum}

\titleformat{name=\chapter,numberless}[block]{\LARGE\bfseries\sffamily\centering}{}{0pt}{\printtitleNoNum}
\newlength\pretitlewidth
\newcommand\bigchapternumber{\resizebox{24pt}{!}{\mdseries\thechapter}}

\newcommand{\printtitleNum}[1]{%
  \settowidth{\pretitlewidth}{{\fontfamily{phv}\selectfont\MakeUppercase{\chaptertitlename}}\space{\bigchapternumber}\space}%
  \parbox[t]{\dimexpr 0.8\textwidth}{%
    \linespread{1}\selectfont
    {\color{violet}\hrule height 4pt}\vskip 0.2ex{\color{violet}\hrule height 1pt}
    \vspace{2ex}\bfseries #1\vspace{2ex}{\color{violet}\hrule height 1pt}}}

\newcommand{\printtitleNoNum}[1]{%
  \parbox[c]{\textwidth}{\linespread{1}\selectfont\centering
    {\color{violet}\hrule height 4pt}\vskip 0.2ex{\color{violet}\hrule height 1pt}
    \vspace{2ex}\bfseries #1\vspace{2ex}{\color{violet}\hrule height 1pt}}}

\titlespacing*{\chapter}{0pt}{-1cm}{1.5cm}

\usepackage{framed}
\definecolor{shadecolor}{rgb}{0.7,0.8,0.8}
\usepackage{longtable,array,multirow,listings,float,subcaption,enumitem,parskip}
\tcbuselibrary{breakable,skins}

\definecolor{primary}{RGB}{0,90,160}
\definecolor{codebg}{RGB}{245,245,250}
\definecolor{codefg}{RGB}{30,30,80}
\definecolor{codegreen}{RGB}{0,150,80}

\lstset{
  language=Python, backgroundcolor=\color{codebg},
  basicstyle=\ttfamily\footnotesize\color{codefg},
  keywordstyle=\bfseries\color{primary},
  commentstyle=\itshape\color{gray},
  stringstyle=\color{codegreen!60!black},
  showstringspaces=false, breaklines=true, frame=single,
  rulecolor=\color{primary!30}, numbers=left,
  numberstyle=\tiny\color{gray}, xleftmargin=1.5em,
  framexleftmargin=1.2em, captionpos=b,
}

\newtcolorbox{infobox}[1][]{colback=primary!5,colframe=primary!60,
  fonttitle=\bfseries\color{primary},title={#1},breakable}

\newtcolorbox{warnbox}[1][]{colback=orange!5,colframe=orange!70,
  fonttitle=\bfseries\color{orange!80!black},title={#1},breakable}

% =============================================================================
\begin{document}
% =============================================================================

% ── Title page ────────────────────────────────────────────────────────────────
\begin{titlepage}
  \newcommand{\bordure}{
    \begin{tikzpicture}[remember picture,overlay]
      \draw[thick]([yshift=-17pt,xshift=17pt]current page.north west)--
        ([yshift=-17pt,xshift=-17pt]current page.north east)--
        ([yshift=17pt,xshift=-17pt]current page.south east)--
        ([yshift=17pt,xshift=17pt]current page.south west)--cycle;
      \draw[ultra thick]([yshift=-15pt,xshift=15pt]current page.north west)--
        ([yshift=-15pt,xshift=-15pt]current page.north east)--
        ([yshift=15pt,xshift=-15pt]current page.south east)--
        ([yshift=15pt,xshift=15pt]current page.south west)--cycle;
      \draw[thick]([yshift=-13pt,xshift=13pt]current page.north west)--
        ([yshift=-13pt,xshift=-13pt]current page.north east)--
        ([yshift=13pt,xshift=-13pt]current page.south east)--
        ([yshift=13pt,xshift=13pt]current page.south west)--cycle;
    \end{tikzpicture}\par\vspace{-2mm}}
  \bordure
  \begin{center}
    \includegraphics[scale=1.3]{aimssn_logo.png}\\[0.7cm]
    \scalebox{1.4}{\color{purple} African Institute for Mathematical Sciences (AIMS), Senegal}\\[0.7cm]
    \scalebox{1.}[1.4]{\fontfamily{phv}\selectfont\textbf{\large\textsc{Computer Vision Project}}}\\[0.3cm]
    \rule{0.5cm}{0.01mm}\textbf{Group 10}\rule{0.5cm}{0.01mm}\\[0.7cm]
    \rule{4.7cm}{1mm}\textsc{\fontfamily{phv}\selectfont\textbf{ Topic }}\rule{4.7cm}{1mm}\\[0.2cm]
    \begin{tcolorbox}[skin=beamer,beamer,shadow={5.5mm}{-3.5mm}{2mm}{fill=blue,opacity=0.25}]
      \begin{minipage}{\linewidth}\begin{center}\vspace{0.2cm}
        \fontfamily{phv}\selectfont\large\bf
        \MakeUppercase{Real-Time Road Traffic Object Detection, Tracking, and Counting}
      \end{center}\end{minipage}
    \end{tcolorbox}\\[1cm]
    \textbf{\fontfamily{phv}\selectfont\uuline{This work was done by:}}\\[0.6cm]
    \begin{minipage}{.8\linewidth}\begin{flushleft}
      1.\ Moute Jean-Baptiste\\[0.3em]
      2.\ Okere Rafiatou\\[0.3em]
      3.\ Gaolatlhe Angelah Kgato
    \end{flushleft}\end{minipage}\\[0.8cm]
    \textbf{\uuline{Supervised by:}}\\[0.25cm]
    \begin{tcolorbox}[rightrule=3mm]\begin{center}
      \scalebox{1.4}{\textbf{Dr. Jordan Felicien Masakuna}}\\[0.3cm]
      {\it\fontfamily{lmss}\selectfont\textcolor{blue!50!black}{University of Kinshasa, Dem. Rep. of Congo.}}\\[0.2cm]
      \underline{Email}: \url{jordan@aims.ac.za}
    \end{center}\end{tcolorbox}\\[0.6cm]
    {\small\today}\\[0.2cm]
    {\scriptsize\it Submitted in partial fulfilment of the requirements for the award of
    Master of Science in Mathematical Sciences at AIMS Senegal}
  \end{center}
\end{titlepage}

% ── Front matter ──────────────────────────────────────────────────────────────
\pagenumbering{roman}\setcounter{page}{1}

\chapter*{Abstract}
\addcontentsline{toc}{chapter}{Abstract}

This report describes a real-time road traffic monitoring system developed as part of the AIMS Senegal Computer Vision Project 2025--2026. The system detects and tracks road-traffic objects — cars, trucks, buses, motorcycles, bicycles, and pedestrians — using a fine-tuned YOLOv11n model combined with the ByteTrack tracking algorithm. Each detected object is assigned a unique identity that persists across frames, allowing the system to count how many distinct objects pass through a scene rather than simply counting detections per frame.

The system was tested on two different traffic scenes and produced a combined total of 397 unique object detections. All results are saved in a standardised 19-column CSV format compatible with the global AIMS traffic dashboard. The application is accessible through a Flask web interface and is deployed publicly on Hugging Face Spaces.

\noindent\textbf{Keywords:} object detection, multi-object tracking, YOLOv11, ByteTrack, traffic monitoring, Flask, fine-tuning.

\tableofcontents\newpage
\pagenumbering{arabic}\setcounter{page}{1}

% =============================================================================
\chapter{Introduction}
% =============================================================================

\section{Motivation}

Road traffic is a daily challenge in many African cities. Infrastructure is often stretched beyond capacity, and city planners lack reliable data to decide where improvements are most needed. Manual counting — sending observers to stand at intersections — is slow, expensive, and human error is unavoidable.

Computer vision provides a better solution. A single camera connected to a detection pipeline can automatically count vehicles and pedestrians, record their types, and build a timestamped record of how traffic evolves through the day. This kind of objective data helps transport authorities identify congestion hotspots, justify infrastructure budgets, and plan road improvements that match actual mobility patterns.

This project is one piece of a larger collaborative effort at AIMS Senegal. Each group processes different traffic videos and contributes their logs to a shared dataset. When all groups' data is combined, the result is a multi-scene traffic monitoring dataset that no single group could have built alone.

\section{Objectives}

\begin{enumerate}[label=\arabic*.]
  \item Detect six types of road-traffic objects using YOLOv11n.
  \item Track each object across frames with a persistent unique ID using ByteTrack.
  \item Count unique objects per scene — not detections, but actual distinct objects.
  \item Log all detections in the shared 19-column CSV schema.
  \item Provide speed estimation, direction detection, and line-crossing events as bonus features.
  \item Deliver a working web application with live stream, counter panel, and analytics dashboard.
\end{enumerate}

% =============================================================================
\chapter{Background}
% =============================================================================

\section{Object Detection with YOLO}

YOLO (You Only Look Once) \cite{redmon2016yolo} processes an entire image in a single forward pass through a neural network and outputs bounding boxes and class probabilities for all objects at once. This makes it much faster than older two-stage detectors that first propose regions and then classify them.

YOLOv11, released by Ultralytics in 2024 \cite{ultralytics2024}, is the latest version of this family. The nano variant (\texttt{yolo11n.pt}) is designed to run on modest hardware, including CPU, while keeping detection quality acceptable. It was originally trained on the COCO dataset, which already includes all six traffic classes we need.

\section{Multi-Object Tracking with ByteTrack}

Detecting objects frame by frame is not enough for counting. If the same car appears in 200 consecutive frames, we need to know it is the same car — not 200 different cars. This is the job of a tracker.

ByteTrack \cite{zhang2022bytetrack} works by associating bounding boxes between frames using the IoU (Intersection over Union) between predicted and detected positions. Its key innovation is that it does not discard low-confidence detections. Instead, it uses them in a second pass to recover tracks that were momentarily lost due to occlusion. This makes it much more robust than simpler trackers like SORT.

ByteTrack is built into Ultralytics YOLO and is activated with a single argument: \texttt{model.track(frame, persist=True)}.

% =============================================================================
\chapter{System Design}
% =============================================================================

\section{Architecture Overview}

The system has three layers:

\begin{enumerate}[label=\arabic*.]
  \item \textbf{Processing layer} — reads video frames, runs YOLOv11 + ByteTrack, computes speed and direction, writes CSV logs.
  \item \textbf{Web layer} — a Flask application with HTTP endpoints, background processing threads, and a JPEG polling stream.
  \item \textbf{Presentation layer} — three HTML pages: landing page, live detection interface, and analytics dashboard.
\end{enumerate}

\section{Detection and Tracking Pipeline}

For each video frame, the \texttt{TrafficTracker.process\_frame()} method:

\begin{enumerate}[label=\arabic*.]
  \item Calls \texttt{model.track(frame, persist=True)} to run detection and tracking together.
  \item For each tracked box, reads the track ID, class, confidence, and bounding-box coordinates.
  \item Computes the centroid $(c_x, c_y)$ and, using the previous centroid $(p_x, p_y)$:
  \begin{align}
    \text{speed} &= \frac{\sqrt{(c_x-p_x)^2+(c_y-p_y)^2}}{\Delta t} \;\text{[px/s]} \\
    \text{direction} &= \begin{cases}\text{up} & c_y < p_y \\ \text{down} & c_y > p_y\end{cases}
  \end{align}
  \item Checks for line crossing: if the centroid moved across $L = \lfloor H/2\rfloor$ between frames.
  \item Writes one row to the CSV log.
  \item Adds the track ID to \texttt{track\_ids\_seen} (a set) if not already present, incrementing the unique count by one.
\end{enumerate}

The unique-counting logic is simple but important: a set guarantees that each track ID is counted exactly once, no matter how many frames it appears in.

\section{CSV Log Schema}

All groups use the same 19-column format (Table~\ref{tab:schema}) so their logs can be merged into the global dataset without transformation.

\begin{table}[H]
  \centering\small
  \caption{Shared 19-column CSV schema.}
  \label{tab:schema}
  \begin{tabular}{lll}
    \toprule
    \textbf{Column} & \textbf{Type} & \textbf{Description}\\
    \midrule
    frame & int & Frame index (0-based)\\
    timestamp\_sec & float & Seconds from video start\\
    scene\_name & string & Scene label\\
    group\_id & string & \texttt{group\_10}\\
    video\_name & string & Video filename\\
    track\_id & int & ByteTrack persistent ID\\
    class\_name & string & car, truck, bus, etc.\\
    confidence & float & Detection score $\in[0,1]$\\
    bbox\_x1/y1/x2/y2 & int & Bounding box corners (px)\\
    cx, cy & int & Centroid coordinates (px)\\
    frame\_width/height & int & Frame dimensions (px)\\
    crossed\_line & bool & Crossed the counting line\\
    direction & string & up or down\\
    speed\_px\_s & float & Speed in pixels per second\\
    \bottomrule
  \end{tabular}
\end{table}

% =============================================================================
\chapter{Implementation}
% =============================================================================

\section{Technology Stack}

\begin{table}[H]
  \centering
  \caption{Tools and libraries used.}
  \begin{tabular}{ll}
    \toprule
    \textbf{Purpose} & \textbf{Tool}\\
    \midrule
    Detection + tracking & Ultralytics YOLOv11 ($\geq$ 8.3)\\
    Video processing & OpenCV headless ($\geq$ 4.9)\\
    Web framework & Flask ($\geq$ 3.0)\\
    YouTube stream & yt-dlp ($\geq$ 2024.1)\\
    Deployment & Docker + Hugging Face Spaces\\
    Frontend & Bootstrap 5, Chart.js\\
    \bottomrule
  \end{tabular}
\end{table}

\section{Fine-Tuning YOLOv11 — Challenges and Lessons Learned}

Fine-tuning a YOLO model means taking pre-trained weights from COCO and continuing training on a new, smaller, task-specific dataset. The goal is to specialise the model for a particular type of scene or object. In practice, this turned out to be the most difficult part of the project.

\subsection{What went wrong}

We encountered two main problems:

\begin{warnbox}[Problem 1 — Training was extremely slow]
Each training run took more than 8 hours for just 10 epochs on a CPU. This made it impossible to experiment with different settings in any reasonable time. The root cause was that we were training on a CPU instead of a GPU. YOLO was designed for GPU training: what takes 8 hours on a CPU takes roughly 15--20 minutes on a mid-range GPU.
\end{warnbox}

\begin{warnbox}[Problem 2 — The fine-tuned model detected fewer objects than the original]
After fine-tuning, the model performed worse than the off-the-shelf COCO weights. Detections were missed, confidence scores dropped, and some classes stopped appearing entirely. This phenomenon is called \textbf{catastrophic forgetting}: when a model is trained on a new small dataset with a learning rate that is too high, it overwrites the features it learned during the original large-scale training.
\end{warnbox}

\subsection{Why fine-tuning can make things worse}

There are three main reasons a fine-tuned model performs worse:

\begin{enumerate}[label=\arabic*.]
  \item \textbf{Dataset too small.} YOLO was trained on 118,000 images. If you fine-tune on a few hundred images, the model overfits — it memorises your small dataset and loses its ability to generalise.
  \item \textbf{Learning rate too high.} A high learning rate destroys the feature representations the model already learned. For fine-tuning, the learning rate should be at least 10$\times$ smaller than the one used for original training.
  \item \textbf{Too many epochs without freezing.} Updating all layers from the start modifies even the low-level features (edges, textures) that do not need to change. Freezing the backbone for the first few epochs and only training the detection head avoids this.
\end{enumerate}

\subsection{The right approach to fine-tuning}

Table~\ref{tab:finetune} summarises the parameters that matter most and the values we recommend based on what we learned.

\begin{table}[H]
  \centering
  \caption{Key fine-tuning parameters and recommended values.}
  \label{tab:finetune}
  \begin{tabular}{lll}
    \toprule
    \textbf{Parameter} & \textbf{What it controls} & \textbf{Recommended value}\\
    \midrule
    \texttt{lr0} & Initial learning rate & 0.0001 -- 0.001 (low)\\
    \texttt{freeze} & Number of layers to freeze & 10 (freeze backbone)\\
    \texttt{epochs} & Training duration & 30--50\\
    \texttt{batch} & Images per step & 8--16 (CPU) / 32+ (GPU)\\
    \texttt{imgsz} & Input resolution & 640\\
    \texttt{patience} & Early stopping rounds & 10\\
    \bottomrule
  \end{tabular}
\end{table}

\begin{lstlisting}[caption=Recommended fine-tuning configuration.]
from ultralytics import YOLO

model = YOLO("yolo11n.pt")
model.train(
    data="traffic.yaml",
    epochs=50,
    imgsz=640,
    batch=8,
    lr0=0.0005,    # low learning rate — preserve pre-trained features
    freeze=10,     # freeze backbone, only train the detection head
    patience=10,   # stop early if no improvement
    device="cpu",
)
\end{lstlisting}

\subsection{When the base model already works well}

YOLOv11n was trained on COCO, which includes all six classes we need (car, truck, bus, motorcycle, bicycle, person). If the base model already detects these objects with acceptable confidence in your videos, fine-tuning may not be necessary and can even hurt performance.

\begin{infobox}[Rule of thumb]
Only fine-tune if your videos contain something the COCO model has not seen before — a very specific camera angle, unusual lighting, or objects that look different from standard COCO images. If the base model already works, use it directly. If you do fine-tune, use a very low learning rate and freeze the backbone.
\end{infobox}

\section{Web Application}

The Flask backend runs video processing in a background thread and shares state with the HTTP server through a thread-safe dictionary. The frontend polls two endpoints every 150 ms and 800 ms respectively:

\begin{itemize}
  \item \texttt{GET /frame} — returns the latest annotated frame as a JPEG image.
  \item \texttt{GET /status} — returns counts, progress, and any error message as JSON.
\end{itemize}

We chose polling over MJPEG streaming because cloud platforms like Hugging Face Spaces use a reverse proxy that terminates long-lived HTTP connections, causing MJPEG streams to go black after a short time. Individual JPEG requests are stateless and work reliably behind any proxy.

\section{Deployment}

The application is containerised with Docker and deployed on Hugging Face Spaces (CPU Basic). The live URL is:
\begin{center}
  \url{https://huggingface.co/spaces/Rafiatou/trafficvision-group10}
\end{center}
The source code is available at:
\begin{center}
  \url{https://github.com/okere-rafiatou/trafficvision-group10}
\end{center}

% =============================================================================
\chapter{Results}
% =============================================================================

\section{Test Scenes}

\begin{table}[H]
  \centering
  \caption{Description of the two processed scenes.}
  \begin{tabular}{llll}
    \toprule
    \textbf{Scene} & \textbf{Resolution} & \textbf{FPS} & \textbf{Type}\\
    \midrule
    Scene 1 (\texttt{traffic.mp4}) & 1280$\times$720 & 30 & Urban intersection\\
    Scene 2 (\texttt{2881980-uhd.mp4}) & 3840$\times$2160 & 24 & Mixed urban road\\
    \bottomrule
  \end{tabular}
\end{table}

\section{Unique Object Counts}

\begin{table}[H]
  \centering
  \caption{Unique objects counted per class and per scene.}
  \label{tab:counts}
  \begin{tabular}{lrrr}
    \toprule
    \textbf{Class} & \textbf{Scene 1} & \textbf{Scene 2} & \textbf{Total}\\
    \midrule
    Car        & 140 & 69  & 209\\
    Truck      & 46  & 3   & 49\\
    Bus        & 5   & 5   & 10\\
    Motorcycle & 1   & 0   & 1\\
    Bicycle    & 0   & 0   & 0\\
    Person     & 0   & 128 & 128\\
    \midrule
    \textbf{Total} & \textbf{192} & \textbf{205} & \textbf{397}\\
    \bottomrule
  \end{tabular}
\end{table}

Cars dominated Scene 1, which is an intersection with dense vehicle flow. Scene 2 had a very high pedestrian count (128 persons), which reflects the more urban character of that footage.

\section{Detection Observations}

\begin{description}
  \item[Confidence scores.] Cars and trucks were detected with confidence typically between 0.78 and 0.93. Motorcycles and bicycles scored lower (0.55--0.72), mainly because they are smaller in the frame at typical camera distances.

  \item[Tracking stability.] ByteTrack maintained stable IDs through most occlusion events. The most common source of ID switches was when two vehicles of the same class overlapped briefly and then separated. This is a known limitation of IoU-based trackers.

  \item[Speed estimation.] Speed is reported in pixels per second. In Scene 1, cars moved at an average of approximately 180--320 px/s depending on their distance from the camera. Without camera calibration, converting to km/h is not possible, but the relative speeds are consistent and meaningful for comparison.

  \item[Direction detection.] Most vehicles in Scene 1 moved downward through the frame (approaching the camera), while pedestrians in Scene 2 showed more varied directions.
\end{description}

\begin{figure}[H]
  \centering
  \fbox{\parbox{0.85\textwidth}{\centering\vspace{3cm}
    [Screenshot: live detection interface with bounding boxes and track IDs]
  \vspace{3cm}}}
  \caption{Live detection interface showing annotated frames, track IDs, and real-time object counters.}
  \label{fig:app}
\end{figure}

\begin{figure}[H]
  \centering
  \fbox{\parbox{0.85\textwidth}{\centering\vspace{3cm}
    [Screenshot: analytics dashboard — class distribution and timeline charts]
  \vspace{3cm}}}
  \caption{Analytics dashboard with per-scene counts, donut chart, and traffic intensity timeline.}
  \label{fig:dashboard}
\end{figure}

% =============================================================================
\chapter{Limitations and Conclusion}
% =============================================================================

\section{Limitations}

\begin{description}
  \item[CPU speed.] YOLOv11n runs at 3--8 frames per second on CPU. The live stream lags behind real time. A GPU would bring this to 30+ fps.
  \item[Speed units.] Speed is in pixels per second. Camera calibration would be needed to convert to km/h.
  \item[Bicycles undetected.] Neither scene contained clearly visible bicycles. The counter remained at zero for this class.
  \item[Fine-tuning results.] Our fine-tuned model performed similarly to or slightly below the base COCO model. Given the small training dataset and the absence of a GPU, the improvement was not significant. The base model is used in the final system.
  \item[YouTube blocking.] Cloud servers are frequently blocked by YouTube's bot detection, making the YouTube input tab unreliable on the hosted version.
\end{description}

\section{Conclusion}

This project built a complete, end-to-end traffic monitoring system — from raw video to structured data to a public web application. Despite the challenges with fine-tuning (slow CPU training, catastrophic forgetting on a small dataset), the system produces reliable unique-object counts and correct CSV logs in the shared format.

The most important lessons were about fine-tuning: the base YOLOv11n model already covers our target classes well because it was trained on COCO. Fine-tuning only helps when the new dataset is large, diverse, and annotated correctly. With a small dataset and no GPU, the base model is the better choice.

Across two scenes the system counted 397 unique objects. All logs follow the shared 19-column schema and are available for integration into the global AIMS traffic dashboard.

% =============================================================================
\begin{thebibliography}{9}

\bibitem{redmon2016yolo}
J. Redmon, S. Divvala, R. Girshick, and A. Farhadi,
``You Only Look Once: Unified, Real-Time Object Detection,''
\textit{CVPR}, 2016.

\bibitem{zhang2022bytetrack}
Y. Zhang et al.,
``ByteTrack: Multi-Object Tracking by Associating Every Detection Box,''
\textit{ECCV}, 2022.

\bibitem{ultralytics2024}
Ultralytics, ``YOLOv11 Documentation,''
\url{https://docs.ultralytics.com}, 2024.

\bibitem{flask}
A. Ronacher, ``Flask: A lightweight WSGI web application framework,''
\url{https://flask.palletsprojects.com}, 2024.

\bibitem{opencv}
G. Bradski, ``The OpenCV Library,''
\textit{Dr. Dobb's Journal of Software Tools}, 2000.

\bibitem{pexels}
Pexels, ``Free Stock Videos,''
\url{https://www.pexels.com/search/videos/traffic/}, 2024.

\end{thebibliography}

\end{document}
